\chapter{Evaluation}\label{evaluation}

Dieses Kapitel evaluiert das in Kapitel~\ref{systemkonzeption} konzipierte und in Kapitel~\ref{implementierung} realisierte System empirisch und liefert damit die Grundlage zur Beantwortung der Forschungsfrage (vgl. Abschnitt~\ref{forschungsfrage}). Es gliedert sich in die Evaluation der KI-Textgenerierung (Abschnitt~\ref{evaluation-der-ki-textgenerierung}), Performance und Laufzeitverhalten (Abschnitt~\ref{performance-und-laufzeitverhalten}), das Latenzpotenzial der Betriebsmodi (Abschnitt~\ref{evaluation-der-betriebsmodi}) sowie den Anforderungsabgleich als Soll-Ist-Vergleich (Abschnitt~\ref{anforderungsabgleich-soll-ist-vergleich}).


\section{Evaluation der
KI-Textgenerierung}\label{evaluation-der-ki-textgenerierung}

Die Evaluation der KI-Textgenerierung bewertet, inwiefern das System journalistische Qualitätsanforderungen hinsichtlich Korrektheit, Tonalität und Vollständigkeit erfüllt. Sie gliedert sich in den Stilvergleich (Abschnitt~\ref{einfluss-der-stilprofile}), die Fehleranalyse (Abschnitt~\ref{qualitative-analyse-der-generierten-texte}), die LLM-as-Judge-Evaluation (Abschnitt~\ref{erweiterte-evaluation}) sowie die funktionale Verifikation der Mehrsprachigkeit (Abschnitt~\ref{mehrsprachige-verifikation}).

\subsection{Stilvergleich}\label{einfluss-der-stilprofile}

Für einen kontrollierten Vergleich wurde dasselbe Event in allen drei Stilen generiert. Als Kontext diente ein Tor von Omar Marmoush in der 78.~Minute (Eintracht Frankfurt vs.~FC Bayern München), womit Stilunterschiede nicht durch unterschiedliche Ereignistypen verfälscht werden.

{\def\LTcaptype{none} % do not increment counter
\begin{longtable}[]{@{}
  >{\raggedright\arraybackslash}p{(\linewidth - 2\tabcolsep) * \real{0.0714}}
  >{\raggedright\arraybackslash}p{(\linewidth - 2\tabcolsep) * \real{0.9286}}@{}}
\toprule\noalign{}
\begin{minipage}[b]{\linewidth}\raggedright
Stil
\end{minipage} & \begin{minipage}[b]{\linewidth}\raggedright
Generierter Text
\end{minipage} \\
\midrule\noalign{}
\endhead
\bottomrule\noalign{}
\endlastfoot
euphorisch & „78. Minute: MARMOOOSH! Der Ball zappelt im Netz!
Unglaublich! Die Hütte explodiert! Was für ein Strahl! Eintracht führt!
TOOOOR!'' \\
neutral & „78. Minute: TOOOOR! Omar Marmoush bringt Eintracht Frankfurt
in Führung!'' \\
kritisch & „78. Minute: TOOOOR! Marmoush mit dem Ausgleich für die
Eintracht! Die Bayern-Abwehr pennt. Frankfurt nutzt die Chance
eiskalt.'' \\
\end{longtable}
}

Die Profile unterscheiden sich klar: \texttt{euphorisch} arbeitet mit Wiederholung, Großschreibung und emotionalen Ausrufen; \texttt{neutral} bleibt faktisch und knapp; \texttt{kritisch} liefert eine Einschätzung statt reiner Beschreibung. Alle drei Ausgaben sind faktenbasiert und genrekonform.

\subsection{Fehleranalyse}\label{qualitative-analyse-der-generierten-texte}

40 KI-generierte Ticker-Einträge aus 15 Bundesliga-Spielen der Saison~2024/25 wurden manuell nach \textbf{Korrektheit}, \textbf{Tonalität} und \textbf{Vollständigkeit} bewertet (Skala 1--5; Bewerter: Entwickler). Der Gesamtdurchschnitt liegt bei 4{,}3. Die Analyse identifizierte folgende Fehlerklassen:

{\footnotesize\def\LTcaptype{none} % do not increment counter
\begin{longtable}[]{@{}
  >{\raggedright\arraybackslash}p{(\linewidth - 6\tabcolsep) * \real{0.2300}}
  >{\raggedright\arraybackslash}p{(\linewidth - 6\tabcolsep) * \real{0.1600}}
  >{\raggedright\arraybackslash}p{(\linewidth - 6\tabcolsep) * \real{0.2900}}
  >{\raggedright\arraybackslash}p{(\linewidth - 6\tabcolsep) * \real{0.3200}}@{}}
\toprule\noalign{}
\begin{minipage}[b]{\linewidth}\raggedright
Fehlerklasse
\end{minipage} & \begin{minipage}[b]{\linewidth}\raggedright
Häufigkeit
\end{minipage} & \begin{minipage}[b]{\linewidth}\raggedright
Beschreibung
\end{minipage} & \begin{minipage}[b]{\linewidth}\raggedright
Beispiel
\end{minipage} \\
\midrule\noalign{}
\endhead
\bottomrule\noalign{}
\endlastfoot
\textbf{Stil-Inkonsistenz} & 8 / 40 (20\,\%) & Neutraler Stil enthält
emotionale Formulierungen & „33. Minute: Oh nein! Lienhart sieht Gelb!
\ldots{} Aber Kopf hoch, Jungs! Kämpfen!'' \emph{(angefordert:
neutral)} \\
\textbf{Fakten-Halluzination} & 2 / 40 (5\,\%) & Das LLM erfindet
Details, die nicht im Kontext stehen & Pre-Match-Text erfand eine
Wettempfehlung: „Wir tippen auf Double Chance: Sieg \ldots{} oder
Remis.'' \\
\textbf{Team-Verwechslung} & 0 / 40 (0\,\%) & Tor wird dem falschen Team
zugeordnet & Nicht aufgetreten \\
\textbf{Wiederholung} & 5 / 40 (13\,\%) & Aufeinanderfolgende Texte
ähneln sich stark im Satzbau & Mehrere Tore mit identischer Formel „X
bringt Y in Führung! Nach Vorlage von Z. 1:0.'' \\
\textbf{Überlänge} & 0 / 40 (0\,\%) & Text überschreitet die typische
Liveticker-Kürze & Nicht aufgetreten \\
\end{longtable}
}

\emph{Beispiel: Tor (euphorisch, Bewertung 5\,/\,5\,/\,5):}

\begin{quote}
„81. Minute: WAS IST DENN HIER LOS?! Collins mit der butterweichen Flanke! DOAN! Der Ball ist drin! TOOOOR! 1:0! Die Hütte bebt!'' \emph{(R. Doan, Eintracht Frankfurt vs.~FSV Mainz 05, Vorlage Collins)}
\end{quote}

Die häufigste Fehlerklasse ist Stil-Inkonsistenz beim neutralen Profil (8/40, 20\,\%); Fakten-Halluzinationen traten mit 5\,\% selten auf.


\subsection{LLM-as-Judge-Evaluation}\label{erweiterte-evaluation}

N\,=\,80 Einträge aus 30 Bundesliga-Spielen der Saison~2025/26 wurden automatisiert von \texttt{claude-sonnet-4-5} nach denselben Dimensionen bewertet wie in Abschnitt~\ref{qualitative-analyse-der-generierten-texte} (Korrektheit, Tonalität, Vollständigkeit, je 1--5); als Generierungsmodell diente \texttt{google/gemini-2.0-flash-lite-001} (LLM-as-Judge-Verfahren nach Zheng et al.~2023).

{\def\LTcaptype{none} % do not increment counter
\begin{longtable}[]{@{}
  >{\raggedright\arraybackslash}p{(\linewidth - 10\tabcolsep) * \real{0.14}}
  >{\raggedright\arraybackslash}p{(\linewidth - 10\tabcolsep) * \real{0.06}}
  >{\raggedright\arraybackslash}p{(\linewidth - 10\tabcolsep) * \real{0.20}}
  >{\raggedright\arraybackslash}p{(\linewidth - 10\tabcolsep) * \real{0.20}}
  >{\raggedright\arraybackslash}p{(\linewidth - 10\tabcolsep) * \real{0.25}}
  >{\raggedright\arraybackslash}p{(\linewidth - 10\tabcolsep) * \real{0.15}}@{}}
\toprule\noalign{}
Stil & n & Korrektheit & Tonalität & Vollständigkeit & Gesamt \\
\midrule\noalign{}
\endhead
\bottomrule\noalign{}
\endlastfoot
neutral     & 27 & 4{,}56 & 4{,}93 & 4{,}78 & 4{,}76 \\
kritisch    & 26 & 4{,}35 & 3{,}50 & 3{,}62 & 3{,}82 \\
euphorisch  & 27 & 3{,}85 & 4{,}41 & 4{,}07 & 4{,}11 \\
\textbf{Gesamt} & \textbf{80} & \textbf{4{,}25} & \textbf{4{,}29} & \textbf{4{,}16} & \textbf{4{,}23} \\
\end{longtable}
}

\emph{K = Korrektheit, T = Tonalität, V = Vollständigkeit (je 1--5). Bewerter: Claude Sonnet~4.5 via OpenRouter. Datenbasis: Football API~v3, Bundesliga Saison 2025/26.}

Das neutrale Stilprofil erzielt den höchsten Gesamtwert (4{,}76) mit besonders starker Tonalität (4{,}93). Das kritische Profil zeigt die größte Schwäche bei Tonalität (3{,}50) und Vollständigkeit (3{,}62), was auf die höheren Anforderungen des analytisch-bewertenden Stils hinweist. Das euphorische Profil liegt mit Ø~4{,}11 im Mittelfeld.

\begin{figure}[H]
\centering
\DiagramQualitaetStilprofil
\caption[Qualitätsbewertung nach Stilprofil (LLM-as-Judge-Evaluation, N\,=\,80)]{Qualitätsbewertung nach Stilprofil: LLM-as-Judge-Evaluation (N~=~80; Bewertungsskala 1--5). Datengrundlage: Football API~v3, Bundesliga Saison 2025/26, 30~Spiele.}
\label{fig:qualitaet-stilprofil}
\end{figure}


\subsection{Funktionale Verifikation der Mehrsprachigkeit}\label{mehrsprachige-verifikation}

Zur Verifikation von F7 (Mehrsprachige Textgenerierung) wurden 21 Einträge eines Pokalspiels (FV Engers~07 vs.\ Eintracht Frankfurt) mit dem Sprachparameter \texttt{es} (Spanisch) generiert. Alle 21 Aufrufe lieferten Status \texttt{ok} bei Modell \texttt{google/gemini-2.0-flash-lite-001} (Temperatur~0{,}5). Die Generierung verteilte sich auf drei Stilprofile (\texttt{neutral}, \texttt{euphorisch}, \texttt{kritisch}) und drei Event-Typen (Tore, Gelbe Karten, Wechsel).

Die Ausgaben zeigen, dass sprachspezifische Ausdruckskonventionen im Prompt-Kontext adaptiert werden. Das deutsche Tor-Format \texttt{TOOOOR!} erscheint als \texttt{GOOOOL!}, die Minutenangabe als \texttt{Minuto~44:} statt \texttt{44.~Minute:}. Zwei Beispieltexte illustrieren die Generierung:

\begin{itemize}\tightlist
\item \textbf{Tor, Stilprofil \texttt{neutral}} (44.~Minute): \emph{„MINUTO 44: ¡¿QUÉ ESTÁ PASANDO?! ¡Bahoya! ¡SÍÍÍ! ¡GOOOOL! ¡Esto es un infierno! ¡0-1!"}
\item \textbf{Gelbe Karte, Stilprofil \texttt{euphorisch}} (23.~Minute): \emph{„Minuto 23: ¡Amarilla para Kristensen! ¡El jugador del Frankfurt derriba a nuestro hombre! ¡Falta para el Engers!"}
\end{itemize}

Der Tor-Eintrag mit Stilprofil \texttt{neutral} weist eine euphorische Tonalität auf, konsistent mit dem in Abschnitt~\ref{einfluss-der-stilprofile} dokumentierten Muster der deutschsprachigen Evaluation. Eine qualitative Bewertung der spanischen Ausgaben liegt nicht vor (vgl.\ Abschnitt~\ref{technische-einordnung}).


\section{Performance und Laufzeitverhalten}\label{performance-und-laufzeitverhalten}\label{llm-latenz}

Dieses Kapitel erfasst die Systemlatenz des produktiven Deployments (Abschnitt~\ref{systemlatenz}), vergleicht die Generierungslatenz der drei architektonisch integrierten LLM-Modelle (Abschnitt~\ref{modellvergleich-latenz}) und leitet daraus das Time-to-Publish-Potenzial der Betriebsmodi ab (Abschnitt~\ref{ttp-abschaetzung}).

\subsection{Systemlatenz}\label{systemlatenz}

Erfasste Messgröße war die \textbf{End-to-End-Latenz} (Client → Backend → OpenRouter → Gemini → Backend → Client) als die für die TTP-Analyse relevante Kenngröße. Gemessen wurden 26 sequenzielle HTTP-Aufrufe an das Render-Deployment (\texttt{/api/v1/ticker/generate/\{event\_id\}}) über 13 Bundesliga-Spiele (Saison 2024/25, Eintracht Frankfurt).

Die 26 Aufrufe liegen zwischen 876\,ms und 4\,226\,ms (Median 3\,338\,ms, Mittelwert 3\,136\,ms, Std.\ 641\,ms). Ein einzelner Ausreißer nach unten (876\,ms) deutet auf einen Gemini-seitigen Cache-Treffer hin; alle übrigen Messungen entfallen auf die Kaltgenerierung (2\,413--4\,226\,ms). Die Fehlerrate beträgt 0\,\% (vgl. Abbildung~\ref{fig:latenzen-histogram}).

\begin{figure}[H]
\centering
\DiagramLatenzenHistogram
\caption[LLM-Latenzmessungen: Render-Deployment (N\,=\,26)]{Verteilung der 26 LLM-Latenzmessungen (Render-Deployment, 13 Bundesliga-Spiele): Kaltgenerierung 2\,413--4\,226\,ms, ein Cache-Treffer bei 876\,ms. Median: 3\,338\,ms, P95: 3\,698\,ms.}
\label{fig:latenzen-histogram}
\end{figure}

\subsection{Modellvergleich}\label{modellvergleich-latenz}

Um die Generierungslatenz der drei architektonisch integrierten Modelle zu vergleichen, wurden je N\,=\,10 sequenzielle Aufrufe direkt gegen die OpenRouter-API gesendet -- mit identischem Prompt (Tor-Event, neutraler Stil, Deutsch, Few-Shot-Referenzen). Diese Messung erfasst die reine LLM-Latenz ohne Backend-Overhead.

{\footnotesize\def\LTcaptype{none}
\begin{longtable}[]{@{}
  >{\raggedright\arraybackslash}p{(\linewidth - 8\tabcolsep) * \real{0.36}}
  >{\raggedright\arraybackslash}p{(\linewidth - 8\tabcolsep) * \real{0.08}}
  >{\raggedright\arraybackslash}p{(\linewidth - 8\tabcolsep) * \real{0.18}}
  >{\raggedright\arraybackslash}p{(\linewidth - 8\tabcolsep) * \real{0.18}}
  >{\raggedright\arraybackslash}p{(\linewidth - 8\tabcolsep) * \real{0.18}}@{}}
\toprule\noalign{}
Modell & N & Median (ms) & P95 (ms) & Std.\ (ms) \\
\midrule\noalign{}
\endhead
\bottomrule\noalign{}
\endlastfoot
\texttt{google/gemini-2.5-flash}        & 10 & \textbf{591} & 1\,459 & 295 \\
\texttt{openai/gpt-4.1-mini}            & 10 & 1\,111        & 2\,759 & 521 \\
\texttt{anthropic/claude-haiku-4-5}     & 10 & 1\,668        & \textbf{1\,825} & \textbf{177} \\
\end{longtable}
}

\emph{Direktmessung via OpenRouter (kein Backend-Overhead). Identischer Prompt je Aufruf. Fehlerrate: 0\,\% bei allen drei Modellen.}

Gemini~2.5~Flash erzielt mit 591\,ms Median die kürzeste Generierungslatenz -- 1,9-fach schneller als GPT-4.1-mini und 2,8-fach schneller als Claude~Haiku~4.5. Claude~Haiku~4.5 weist mit einer Standardabweichung von 177\,ms das konstanteste Antwortverhalten auf; GPT-4.1-mini zeigt mit 521\,ms die höchste Streuung. Alle drei Modelle erzielen eine Fehlerrate von 0\,\%.

\subsection{TTP-Abschätzung}\label{ttp-abschaetzung}

Die drei Betriebsmodi unterscheiden sich strukturell in ihrer erreichbaren Publikationslatenz. Auf Basis der gemessenen LLM-Latenzen (Median 3\,338~ms, P95 3\,698~ms, vgl. Abschnitt~\ref{systemlatenz}) und des 15-Sekunden-Polling-Intervalls im Live-Modus (vgl. Abschnitt~\ref{frontend-polling}) ergibt sich für den \texttt{auto}-Modus eine architekturbasierte Time-to-Publish von $\approx$\,15--40~s. Die folgende Tabelle schlüsselt die Latenzkomponenten auf:

{\def\LTcaptype{none}
\begin{longtable}[]{@{}llll@{}}
\toprule\noalign{}
Komponente & Gemessen? & Typischer Wert & Quelle \\
\midrule\noalign{}
\endhead
\bottomrule\noalign{}
\endlastfoot
n8n Datenimport (08\_import\_events) & Ja & 8\,047~ms Median (P95: 18\,269~ms) & N\,=\,82, n8n-Execution-Log \\
n8n Live-Stats-Monitor & Ja & 1\,517~ms Median (P95: 6\,559~ms) & N\,=\,37, n8n-Execution-Log \\
n8n → Backend (HTTP) & Teilw. & $\sim$1\,300~ms & gemessen (äquiv. Netzwerkposition) \\
Backend → LLM → Backend & Ja & 3\,338~ms Median & Abschnitt~\ref{systemlatenz} \\
Backend → DB & Nein & $<$50~ms & geschätzt \\
Frontend-Polling & Ja & bis 15~s & Abschnitt~\ref{frontend-polling} \\
\textbf{Gesamt TTP (\texttt{auto})} & Teilw. & $\approx$\,15--40~s & Potenzialschätzung \\
\end{longtable}
}

Die n8n-Execution-Logs liefern erstmals gemessene Laufzeiten für zwei Workflow-Gruppen: Der Event-Import-Workflow (\texttt{08\_import\_events}, N\,=\,82) erreicht einen Median von 8\,047~ms -- er beinhaltet den Football-API-Aufruf sowie die vollständige Datenbankpersistenz und tritt periodisch auf. Der Live-Stats-Monitor (N\,=\,37) ist mit 1\,517~ms Median deutlich schneller, da er nur Spielstatistiken abgleicht. Zusammen mit der gemessenen Backend-LLM-Latenz (3\,338~ms Median) und dem Frontend-Polling (bis 15~s) ergibt sich eine architekturbasierte TTP von $\approx$\,15--40~s für den \texttt{auto}-Modus. Da ein kontrolliertes Feldexperiment nicht durchführbar war, bleibt dies eine Potenzialabschätzung. Die manuelle Baseline liegt bei 30--120~s (vgl. Abschnitt~\ref{zeitdruck-und-fehleranfaelligkeit}); der Interviewpartner nannte für typische Ereignisse 30--90~s (vgl. Abschnitt~\ref{experteninterview-auswertung}). Der \texttt{coop}-Modus durchläuft dieselbe Pipeline und erreicht eine vergleichbare Generierungslatenz; die Gesamtpublikationslatenz erhöht sich um die variable Redaktionszeit.


\section{Anforderungsabgleich
(Soll-Ist-Vergleich)}\label{anforderungsabgleich-soll-ist-vergleich}\label{technische-qualitätssicherung}

23 der 24 in Abschnitt~\ref{zusammenfassung-der-systemanforderungen} hergeleiteten Anforderungen sind vollständig erfüllt. F7 ist funktional implementiert, eine qualitative Evaluation der Übersetzungsqualität steht als Folgestudie aus (vgl. Abschnitt~\ref{technische-einordnung}). Die folgende Tabelle dokumentiert den Nachweis je Anforderung:

{\def\LTcaptype{none} % do not increment counter
\begin{longtable}[]{@{}
  >{\raggedright\arraybackslash}p{(\linewidth - 4\tabcolsep) * \real{0.06}}
  >{\raggedright\arraybackslash}p{(\linewidth - 4\tabcolsep) * \real{0.38}}
  >{\raggedright\arraybackslash}p{(\linewidth - 4\tabcolsep) * \real{0.56}}@{}}
\toprule\noalign{}
Nr. & Anforderung & Nachweis \\
\midrule\noalign{}
\endhead
\bottomrule\noalign{}
\endlastfoot
\multicolumn{3}{l}{\emph{Funktionale Anforderungen}} \\
F1 & Drei Betriebsmodi (auto, coop, manual) & \texttt{PATCH /matches/\{id\}/ticker-mode}; 5 API-Tests \\
F2 & KI-Generierung für alle Event-Typen & 24 Typen in \texttt{EVENT\_TYPE\_LABEL}; Bulk-Endpoint \\
F3 & Drei Stilprofile & \texttt{STYLE\_DESC} im Prompt; per Match konfigurierbar \\
F4 & Stilkonfiguration über vereinsspezifische Referenztexte & Bis zu 3 Referenzen aus \texttt{style\_references}; nach Event-Typ gefiltert \\
F5 & Ticker-Lifecycle (draft → published / rejected) & State-Machine mit publish/reject/retract; 21 API-Tests \\
F6 & Slash-Command-Parser & 11 Phasen- und 14 Event-Commands (in \texttt{CMD\_MAP}); 45 Unit-Tests \\
F7 & Mehrsprachige Textgenerierung & \texttt{translate\_text()} mit Temperatur 0{,}1; Batch-Endpoint; funktional verifiziert; Übersetzungsqualität qualitativ nicht benchmarkt (vgl. Abschnitt~\ref{technische-einordnung}) \\
F8 & Idempotenter Datenimport & UPSERT-Strategien für Events, Spieler, Teams \\
F9 & Anbieterunabhängige LLM-Integration mit Ausfallsicherheit & 5 Provider (openrouter → gemini → openai → anthropic → mock) \\
F10 & Deduplizierung & Prüfung auf bestehenden Eintrag per \texttt{event\_id} vor LLM-Aufruf \\
F11 & Pre-Match-Kontextgenerierung & 7 Context-Builder in \texttt{llm\_context\_builders.py} \\
F12 & Live-Statistik-Updates & \texttt{ctx\_live\_stats()} in \texttt{llm\_context\_builders.py}; automatische Zwischenstand-Texte; durch \texttt{test\_ticker\_service.py} (Context-Builder-Tests) abgedeckt \\
F13 & Medienintegration (ScorePlay, Social-Media-Panels) & MediaQueue-Architektur; n8n-Workflows 08, 10--12 (ScorePlay, X, YouTube, Instagram); instanzspezifische Panel-Einblendung im Frontend \\
\midrule\noalign{}
\multicolumn{3}{l}{\emph{Nicht-funktionale Anforderungen}} \\
N1 & Concurrency-Begrenzung LLM & \texttt{asyncio.Semaphore(8)} in \texttt{ticker\_service.py} \\
N2 & Retry mit Rate-Limit-Erkennung & 3 Versuche; 30\,s/60\,s bei Rate-Limit; 1\,s/2\,s sonst \\
N3 & Reproduzierbare, isolierte Testausführung unter Produktionsbedingungen & Rollback-Fixtures; keine persistenten Testdaten \\
N4 & Typsichere Konfigurationsschnittstellen zwischen Komponenten & 98{,}02\,\% type-coverage (\texttt{npx type-coverage}); 0 Compiler-Fehler \\
N5 & Responsive UI & Playwright-Test 375×812 (iPhone~X\,/\,XS\,/\,12\,mini u.\,a.); Mobile Tab Bar \\
N6 & Fehlerresistenz Frontend & \texttt{ErrorBoundary} mit Fallback-UI; 4 Tests \\
\midrule\noalign{}
\multicolumn{3}{l}{\emph{Architektur-Anforderungen}} \\
A1 & Schichtenbasierte Systemarchitektur mit klar getrennten Verantwortlichkeiten & n8n (Daten), FastAPI/PostgreSQL (Anwendung), React (Präsentation) \\
A2 & Repository-Pattern & 13 Repository-Klassen \\
A3 & 71 API-Endpunkte unter \texttt{/api/v1} & 13 Router-Dateien; 95 API-Integrationstests \\
A4 & 17 ORM-Modelle / 18 Datenbanktabellen & 100\,\% Model-Coverage \\
A5 & White-Label-Fähigkeit & \texttt{src/config/whitelabel.ts} (Frontend-Branding), \texttt{DEFAULT\_TICKER\_INSTANCE} in \texttt{constants.py} (Backend-Instanz), instanzspezifische Few-Shot-Filterung über \texttt{style\_references} (vgl. Abschnitt~\ref{white-label-konfiguration-impl}) \\
\end{longtable}
}

Die Evaluation liefert die empirische Grundlage für die kritische Einordnung in Kapitel~\ref{diskussion} und die Beantwortung der Forschungsfrage in Abschnitt~\ref{beantwortung-der-forschungsfrage}.
